\subsection{归纳集}

定义3. 若有序集E的每个全序子集在E中都有上界，则称E为归纳的。

\minorheading{示例}

\begin{enumerate}
\def\labelenumi{(\arabic{enumi})}
\item
  令 \(\mathfrak{F}\) 为一个集合 \(A\) 的子集族，按包含关系排序，并且对于每个全序子集 \(\mathfrak{G}\) （ \(\mathfrak{F}\) 的子集）， \(\mathfrak{G}\) 中这些集合的并集属于 \(\mathfrak{F}\) 。则 \(\mathfrak{F}\) 关于关系 \(\subset\) 是归纳的，因为 \(\mathfrak{G}\) 中这些集合的并集是 \(\mathfrak{G}\) 在 \(\mathfrak{P}(A)\) 中的最小上界.
\item
  子集族的一个重要的归纳例子（关于关系 \(\subset\) ）是集合 \(\mathfrak{F}\) （从集合A的子集到集合B的映射的图所成的集合）。因为 \(\mathfrak{F}\) 是 \(\mathfrak{P}(A \times B)\) 的一个子集，并且说一个子集 \(\mathfrak{G}\) （ \(\mathfrak{F}\) 的子集）按包含关系全序，意味着 \(\mathfrak{G}\) 的元素是一些映射的图，使得给定其中任意两个映射，一个是另一个的扩张。由此立即得出， \(\mathfrak{G}\) 中这些集合的并集是 \(\mathfrak{F}\) 的元素 (第二章，§4，第6号，命题7)。因此集合 \(\Phi(A, B)\) （将A的子集映入B的映射族）关于序关系“ \(v\) 是 \(u\) 的延拓”（在 \(u\) 与 \(v\) 之间）是归纳的.
\end{enumerate}

* (3) 由无穷公理（§ 6，第1号）可知，自然整数的良序集关于关系≤不是归纳的。 *

定理2（“佐恩引理”）。每个归纳有序集都有一个极大元。

该定理是以下结果的一个特例：

命题4. 设E是一个有序集，其中每个良序子集都有上界；则E有一个极大元。

我们称一个元素 \(v \in E\) 是某个子集 \(X\) （ \(E\) 的子集）的一个严格上界，若 \(v\) 是 \(X\) 的一个上界并且 \(v \notin X\) . 设 \(\mathfrak{S}\) 为 \(E\) 的这样一些子集——它们在 \(E\) 中有严格上界——所成的集合，并且对于每个 \(S \in \mathfrak{S}\) ，置 \(p(S) = \tau_v(v\) 是 \(S)\) 的严格上界 ；则 \(p(S)\) 是 \(S\) 的一个严格上界。将第3号的引理3应用于 \(\mathfrak{S}\) 与 \(p\) ，我们看到存在一个子集 \(M\) （ \(E\) 的子集）以及一个良序 \(\Gamma\) （在 \(M\) 上的良序），满足引理3的条件；特别地，M在E中没有严格上界。此外，该序关系 \(\Gamma\) 与E上的序在M上诱导出的序完全相同。因为在M中，关系“ \(y \in \Gamma\langle x \rangle\) 并且 \(x \neq y\) ”等价于 \(x \in S_y\) ；并且由于 \(p(S_y) = y\) 是 \(S_y\) 的一个上界（关于E上的序），这意味着在E中 \(x < y\) 。而这意味着M到E的注入是一个严格递增的映射（当M被赋予该序关系 \(\Gamma\) 时）；并且由于 M 是全序的，因此关系 \(y \in \Gamma\langle x \rangle\) 并且 \(x \leqslant y\) 在 M 中等价（ \(\S 1\) ，第 12 号，命题 11）。既然如此，根据假设，存在 M 在 E 中的上界 \(m\) ；但由于 M 没有严格上界，因此得出 \(m\) 是 E 的一个极大元。

推论 1. 设 E 是一个归纳有序集，且 a 是 E 的一个元素。则存在 E 的一个极大元 m，使得 \(m \geqslant a\) .

由此根据定义3可知，E 中满足 \(x \geqslant a\) 的元素所成的集合 F 在E中是归纳的，且 F 的一个极大元在 E 中也是极大元。

推论2。设 \(\mathfrak{F}\) 是集合 E 的一族子集，使得对于每一个按包含关系全序的子集 \(\mathfrak{G}\) （ \(\mathfrak{F}\) 的子集）， \(\mathfrak{G}\) 中集合的并集（相应地，交集）属于 \(\mathfrak{F}\) ；则 \(\mathfrak{F}\) 具有一个极大（相应地，极小）元素。

